% ==============================================================================
% CHƯƠNG 3: DỮ LIỆU VÀ PHƯƠNG PHÁP
% ==============================================================================
\chapter{Dữ liệu và phương pháp}
\label{chap:phuong-phap}

\section{Nguồn dữ liệu}
\label{sec:du-lieu}

Mô tả nguồn, giai đoạn, quy mô và các trường thông tin của dữ liệu.

\section{Quy trình thực hiện}
\label{sec:quy-trinh}

\Hinh{fig:quy-trinh} tóm tắt các bước thực hiện.

\begin{figure}[H]
  \centering
  \begin{tikzpicture}[node distance=1.1cm,
      buoc/.style={rectangle, rounded corners=2mm, draw=bkxanh, thick, fill=bkxanhnen,
                   text width=2.7cm, minimum height=1.1cm, align=center, font=\small\bfseries},
      mui/.style={-Latex, thick, bkxanhsang}]
    \node[buoc] (a) {Thu thập\\dữ liệu};
    \node[buoc, right=of a] (b) {Làm sạch \&\\chuẩn hóa};
    \node[buoc, right=of b] (c) {Phân tích\\mô hình};
    \node[buoc, below=0.9cm of c, fill=xanhlanhat, draw=xanhla] (d) {Đánh giá\\kết quả};
    \draw[mui] (a) -- (b);
    \draw[mui] (b) -- (c);
    \draw[mui] (c) -- (d);
    \draw[mui, dashed] (d) -| node[pos=0.25, below, font=\footnotesize\itshape]{điều chỉnh} (b);
  \end{tikzpicture}
  \caption{Quy trình thực hiện của báo cáo}
  \label{fig:quy-trinh}
\end{figure}

\section{Công cụ}
\label{sec:cong-cu}

Đoạn mã dưới đây minh họa bước làm sạch dữ liệu (xem thêm \PhuLuc{app:du-lieu}).

\begin{khungma}{lam_sach_du_lieu.py}{Python}
import pandas as pd

df = pd.read_csv("tin_tuyen_dung.csv")
df = df.dropna(subset=["ma_nghe"])       # bỏ tin thiếu mã nghề
df["nam"] = pd.to_datetime(df["ngay"]).dt.year
print(df.groupby("nam").size())
\end{khungma}

Môi trường chạy được cài đặt bằng các lệnh sau:

\begin{khungbash}[Cài đặt môi trường]
$ python -m venv .venv
$ source .venv/bin/activate
$ pip install pandas matplotlib   # thư viện xử lý & vẽ biểu đồ
\end{khungbash}

\begin{goiy}[Mẹo trình bày]
  Chỉ đưa vào chương này những đoạn mã ngắn, minh họa ý chính; mã đầy đủ để ở phụ lục.
\end{goiy}
